\chapter{Special Topics}\label{ch:special}

\section[Randomized algorithms]{Randomized algorithms \hfill \plat}

Randomness gives simple algorithms with guarantees that hold with overwhelming probability, and it defeats inputs built to
break deterministic choices.

\begin{itemize}
  \item \textbf{Random pivots and shuffles}: shuffling the input before inserting into a naive BST, before a quickselect, or before
        an incremental algorithm (smallest enclosing circle, expected $\Oh{n}$) turns worst cases into expected cases.
  \item \textbf{Randomized hashing}: random bases and seeds (string hashing, \texttt{custom\_hash}, set hashing by random 64-bit
        values) make collisions both rare and impossible to construct in advance.
  \item \textbf{Sampling}: if an answer is shared by at least half of the elements (a majority, a gcd that divides at least half
        of the numbers, a line through at least a fraction of the points), pick a random element and test it: each try succeeds
        with probability $\ge \frac12$, so 30 tries fail with probability $2^{-30}$.
  \item \textbf{Random priorities}: treaps, random weights in skip lists.
\end{itemize}

\begin{problem}{Ghd (Codeforces 364D)}
Find the largest $g$ that divides at least half of $n \le 10^6$ numbers up to $10^{12}$.
\solution The answer divides some element of the ``good'' half, so a random element is a candidate source with probability
$\ge \frac12$. For a random $a_i$, the answer is a divisor of $a_i$; count, for each divisor of $a_i$, how many elements share it
(via gcds with $a_i$, grouped). A dozen random picks suffice.
\end{problem}

\begin{pitfall}[Seeds]
Seed with the clock (\cc{chrono::steady_clock::now().time_since_epoch().count()}), not a constant. On Codeforces, a fixed seed
lets anyone read your code and craft an anti-random input during the hacking phase.
\end{pitfall}

\usetemplate{Random generator}{Ctrl+Alt+T, Shift+R}{cprng}

\section[Interactive problems]{Interactive problems \hfill \gold}

The judge answers your questions during the run. Rules that always apply:
\begin{itemize}
  \item \textbf{Flush after every query}: \cc{cout << ... << endl} (\cc{endl} flushes) or \cc{cout.flush()}. Forgetting it gives
        ``Idleness limit exceeded''.
  \item \textbf{Count your queries} against the limit before coding; $\log_2(10^9) \approx 30$, $\lceil\log_3 n\rceil$ for ternary
        splits.
  \item \textbf{Stop immediately} if the judge answers $-1$ (wrong query or limit exceeded), or you get a misleading verdict.
  \item Some judges are \textbf{adaptive}: the hidden answer is only fixed to be consistent with your questions, so randomised
        strategies that rely on a fixed hidden value can fail. Design strategies that are worst-case optimal.
\end{itemize}

\usetemplate{Interactive helpers}{Ctrl+Alt+T, Shift+I}{cpinteractive}

\begin{tip}[Testing interactive solutions locally]
CPH cannot play the judge. Write a small interactor that reads your queries and answers them, or temporarily replace \cc{ask}
with a function that computes the answer from a hidden array you generate; the rest of the solution stays unchanged.
\end{tip}

\section[Communication problems]{Communication problems \hfill \ioi}

Your program runs twice: first to \emph{encode} some information into a limited message, then (without memory of the first
run) to \emph{decode} it. Typical tools: fixing a canonical order so both runs agree, encoding numbers in the ``shape'' of the
output (which elements are chosen, their order: $n!$ permutations carry $\log_2 n!$ bits), and error-correcting parity bits. Think
of the two runs as two parties who agreed on a protocol in advance and design the protocol first.

\section[Bitsets]{Bitsets \hfill \plat}

\cc{bitset<N>} packs booleans 64 per machine word, and its operations work word by word. Any algorithm dominated by OR/AND of
boolean rows gets a 64$\times$ speed-up:

\begin{itemize}
  \item subset sums and knapsack feasibility: \cc{b |= b << w};
  \item transitive closure / reachability in a DAG: \cc{reach[u] |= reach[v]} in reverse topological order, $\Oh{nm/64}$;
  \item counting triangles or common neighbours: \cc{(adj[u] & adj[v]).count()};
  \item Gaussian elimination over GF(2) with rows as bitsets.
\end{itemize}

\begin{pitfall}[Bitset sizes are compile-time]
\cc{N} must be a constant. If the needed size varies between tests, use the maximum; or write a template function
instantiated for a few sizes and dispatch at run time.
\end{pitfall}

\section[Matroids]{Matroids \hfill \ioi}

A \defn{matroid} is a ground set with a family of ``independent'' subsets such that (1) subsets of independent sets are
independent and (2) if $A$ and $B$ are independent with $|A| < |B|$, some element of $B \setminus A$ can be added to $A$ keeping it
independent (the \emph{exchange property}). The theorem that makes them useful: \textbf{on a matroid, the greedy algorithm (add
elements in order of decreasing weight whenever independence is kept) finds a maximum-weight independent set.}

\begin{center}
\begin{tabularx}{\linewidth}{p{3.6cm}L}
\toprule
\thd{Matroid} & \thd{Independent sets, and the greedy algorithm it justifies} \\
\midrule
graphic & forests of a graph: Kruskal's MST \\
linear (over GF(2)) & linearly independent vectors: XOR basis with maximum total weight \\
uniform & sets of size $\le k$: take the $k$ heaviest \\
partition & at most $c_i$ elements from group $i$ \\
scheduling & sets of unit jobs that can all meet their deadlines: the deadline greedy of Chapter~\ref{ch:sorting} \\
\bottomrule
\end{tabularx}
\end{center}

\begin{recognize}
A greedy ``sort by weight, take if still feasible'' that you cannot prove with an exchange argument: check whether feasibility
defines a matroid; if it does, greedy is correct. If feasibility is the intersection of \emph{two} matroids (for example ``a
spanning forest using at most one edge of each colour'': graphic $\cap$ partition), greedy fails but \textbf{matroid
intersection} finds the largest common independent set with repeated augmenting paths in an exchange graph, in
$\Oh{r\cdot n\cdot(\text{oracle cost})}$ for rank $r$. Bipartite matching is the simplest case (two partition matroids).
\end{recognize}

\section[Vectorization and pragmas]{Vectorization and pragmas \hfill \ioi}

Modern CPUs process 4--8 integers per instruction (SIMD). GCC vectorizes simple loops over contiguous arrays automatically at
\texttt{-O3}; some contestants add
\begin{code}
#pragma GCC optimize("O3,unroll-loops")
#pragma GCC target("avx2,bmi,bmi2,popcnt,lzcnt")
\end{code}
at the top of the file. It can make an $\Oh{n^2}$ loop over plain arrays several times faster, which occasionally turns an
intended-to-fail $n^2$ into a pass. It is not a substitute for a better algorithm, it helps only branch-free loops over arrays
(not \cc{set}, \cc{map} or recursion), and \texttt{target} can crash with ``illegal instruction'' on judges whose CPUs lack AVX2
(some older USACO and ICPC machines). Keep it out of your default template and add it only when a constant factor is the
bottleneck.

\section[Constructive problems]{Constructive problems \hfill \silver}

``Output any array / permutation / graph with property P'' problems reward patterns more than algorithms:

\begin{itemize}
  \item Solve small $n$ by brute force (write the brute force!), print all solutions, and look for a structure that
        generalises: alternating, sorted then reversed, pairs swapped, a spiral.
  \item Build inductively: from a solution for $n$, get one for $n + 1$ or $n + 2$ by a local change.
  \item Look for invariants that rule out some $n$ (parity, sums modulo something); the impossible cases often show the
        construction for the rest.
  \item Extremal choices: put the largest element first, connect everything to one vertex, use a star or a path.
\end{itemize}
